\section{The actual class-rank identity}
\label{sec:class-rank}

\status{CR1--CR6 prove the literal \texttt{locPack\_d\_eq} and
\texttt{pack\_d\_eq} conclusions from supplied \texttt{D7},
\texttt{LevelData}, \texttt{Forms0}, and the local package of LI.
The class correspondence and the trivial-stabilizer argument are proved
below. The construction of those global records remains a separate task.
An equality of ranks does not establish an integral or metric comparison.}

\paragraph{The quantity to be compared.}
Let $k$ be the coefficient rank and put $m_{\mathrm{loc}}=k-1$;
this local polynomial parameter is distinct from the global chain
length $D.R.m$. Let $\mathcal O\subset K$ be the coefficient DVR and
its fraction field. At a prime $p$, the source has places $v$ and
Frobenius positions $s=0,\ldots,n_v-1$, with a decreasing integer
weight $\lambda^0_{v,s}$ of length $k$. A tuple $\mathbf L=(L_v)_v$
of full $W_v$-lattices gives a polynomial lattice
\[
 H(\mathbf L)=\bigotimes_{v,s}
       D_{\lambda_{v,s}}(\Lambda_{v,s}),\qquad
 \Lambda_{v,s}=\mathcal O\,\iota_v(\sigma_v^s L_v),
 \quad
 \lambda_{v,s}=
 \begin{cases}
 \lambda^0_{v,s},&s\text{ even},\\
 -\operatorname{rev}\lambda^0_{v,s},&s\text{ odd}.
 \end{cases}                                                    \tag{CR.1}
\]
The tensor is realized as products in disjoint polynomial-variable
blocks; it is the literal \texttt{HtupW}.
For each away component $i$, a group $\Gamma_i$ acts on lattice tuples
and polynomial fibers, with an $\mathcal O^{\times}$-valued character
$\chi_i$. The module $M$ consists of functions satisfying
\[
 f_i(\gamma\mathbf L)=\chi_i(\gamma)\rho_i(\gamma)f_i(\mathbf L),
 \qquad f_i(\mathbf L)\in H(\mathbf L).
                                                               \tag{CR.2}
\]
As in the source, these functions are extended by zero to tuples which
are not full lattices. Its rank $d=\operatorname{rank}_{\mathcal O}M$
is the number appearing in the local conclusions.

\paragraph{CR1. The literal Weyl fibers have the required rank.}
Here $D_\lambda$ uses the integral \emph{Weyl} lattice $\Delta$,
not the dual-Weyl coefficient lattice $\nabla$ of CB2.
For $\nu_j=\lambda_j-\lambda_k$, let $\Delta_\nu(\mathbb Z)$ be the
integer span of the divided-power words on its highest polynomial.
All these words have bounded multidegree. Thus this is a finite free
subgroup of a finite free group of integer polynomials. Choose any
integer basis $b_1,\ldots,b_q$. The characteristic-zero argument of IW3
identifies its rational span with that of $\nabla_\nu$; CB2--CB3 give
\[
 q=D_k(\lambda):=
       \prod_{a<b}\frac{\lambda_a-\lambda_b+b-a}{b-a}.
                                                               \tag{CR.3}
\]
Divided-power polynomial operators have integer coefficients and commute
with coefficient extension. Every generator over $\mathcal O$ is
therefore an $\mathcal O$-linear combination of the same $b_j$.
A nonzero integer minor of their coefficient matrix remains nonzero in
$K$, so these images are still independent. This proves the exact
base-change assertion for $\Delta$, without asserting
$\Delta=\nabla$ integrally.

For $\Lambda=g\mathcal O^k$, the source's \texttt{Dl0} is
\[
 D_\lambda(\Lambda)
   =\det(g)^{-\lambda_k}R(g^{-1})\Delta_\nu(\mathcal O),
       \qquad R(g)P(X)=P(gX).                            \tag{CR.4}
\]
This expression is independent of the chosen basis $g$:
changing it to $gh$, $h\in\mathrm{GL}_k(\mathcal O)$, preserves
the Weyl lattice and introduces only an integral determinant unit.
The displayed invertible $K$-linear transformation sends the actual
integer basis to an $\mathcal O$-basis of (CR.4).
Products of these bases in disjoint blocks prove
\[
 \operatorname{rank}_{\mathcal O}H(\mathbf L)
    =q_0:=\prod_{v,s}D_k(\lambda^0_{v,s}).                \tag{CR.5}
\]
Indeed, constant weight shifts and negative reversal preserve the
product in (CR.3). Neither the basis matrices nor their determinant
twists need be integral units.

The required characteristic zero follows from the actual Hida setting.
If $W_v$ had prime characteristic $\ell$, then $p=0$ when $\ell=p$,
contradicting that $p$ is a uniformizer; if $\ell\ne p$, then $p$ is a
unit in the prime subfield, again a contradiction. Thus $W_v$, its
fraction field, and its embedded extension $K$ have characteristic zero.

\paragraph{CR2. Evaluate on one representative of each orbit.}
Suppose $\Gamma_i$ acts freely on the full lattice tuples and has finitely
many orbits. Choose one representative $\mathbf L_x$ of each orbit.
The determinant-twisted row action in (CR.4) takes
$H(\mathbf L)$ onto $H(\gamma\mathbf L)$: apply the formula to a
basis matrix and then apply $\gamma^{-1}$ for the reverse inclusion.
The scalar $\chi_i(\gamma)$ is a unit. Consequently evaluation is an
\emph{integral} linear isomorphism
\[
 M_i\ \longrightarrow\ \prod_{x\in\Gamma_i\backslash\mathrm{Lat}}
                                  H(\mathbf L_x).       \tag{CR.6}
\]
For its inverse, prescribe $a_x\in H(\mathbf L_x)$ and set
$f_i(\gamma\mathbf L_x)=\chi_i(\gamma)\rho_i(\gamma)a_x$.
The orbit determines $x$ and freeness determines $\gamma$, so this
is well-defined. The action and character laws give (CR.2), and the
fiber transport gives its integrality. Conversely, (CR.2) determines
every value from the representatives. This proof also explains why
nontrivial stabilizers would require invariant subspaces rather than
arbitrary fibers.

\paragraph{CR3. The actual class labels are the local lattice orbits.}
Let $\mathcal G=\mathrm U(J)(\mathbb A_{F,f})$, let
$\mathcal H=\mathrm U(J)(F)$ denote its rational subgroup, and let
$U$ be the supplied level. Remove its conditions above $p$ to obtain
$U'$. The local identifications of LI supply a projection and section
\[
 \pi:\mathcal G\longrightarrow
       G_p:=\prod_{v\mid p}\mathrm{GL}_k(Q_v),\qquad
 j:G_p\longrightarrow\mathcal G,
 \qquad \pi j=\operatorname{id}_{G_p},\quad j(G_p)\subset U'.
\]
Here $j$ inserts a split unitary element above $p$ and the identity
elsewhere; its second split component is the adjoint inverse of the
first. In the frames of LI5,
\[
 U=\{a\in U':\pi(a)\in K_p\},
       \qquad K_p=\prod_{v\mid p}\mathrm{GL}_k(W_v).
                                                               \tag{CR.7}
\]
To check this equality for the actual level, $\texttt{GIdx}$ excludes
$p$ from the two auxiliary primes and the integer $T$ controlling
$\Lambda\subset\Lambda^{\#}\subset T^{-1}\Lambda$.
Localization therefore makes $\Lambda$ self-dual. In the integral
splitting $E\otimes_F F_v=F_v\times F_v$, its two projected lattices
are mutually dual. A unitary element preserves both exactly when its
first component preserves the first projected lattice. LI5 identifies
that condition with integrality of the framed matrix and its inverse,
which is precisely membership in $\mathrm{GL}_k(W_v)$.

Choose representatives $y_i$ for $\mathcal H\backslash\mathcal G/U'$
with $\pi(y_i)=1$. They can always be so normalized: replace any
representative $y$ by $yj(\pi(y))^{-1}$. Define
\[
 \Gamma_i=\{\gamma\in\mathcal H:y_i^{-1}\gamma y_i\in U'\}.
\]
Every full lattice tuple has the form $g\mathbf W$, where
$\mathbf W=(W_v^k)_v$, and two matrices give the same tuple exactly
when they differ on the right by $K_p$. Define its class by
\[
 (i,g\mathbf W)\longmapsto[\,y_i j(g)\,]
                        \in\mathcal H\backslash\mathcal G/U.
                                                               \tag{CR.8}
\]
Changing $g$ by $gk$, $k\in K_p$, does not change this class, by
(CR.7). More precisely, multiplication in the double quotient gives
\[
 [y_i j(g)]=[y_{i'}j(g')]
 \quad\Longleftrightarrow\quad
 i=i'\ \text{ and }\
 \exists\gamma\in\Gamma_i:\quad
       g'\mathbf W=\pi(\gamma)g\mathbf W.               \tag{CR.9}
\]
For the reverse implication the needed right-level element is
$j(g)^{-1}y_i^{-1}\gamma^{-1}y_i j(g')$; its membership follows
from (CR.7) and
$g'^{-1}\pi(\gamma)g\in K_p$.
For the forward implication, first pass to the $U'$-quotient to get
$i=i'$, and then project the equality to $G_p$.

There are no missing classes. If $x=\gamma y_i u'$ with $u'\in U'$,
then $j(\pi(u'))^{-1}u'\in U$ by (CR.7). Thus $[x]$ is the class
of $y_i j(\pi(u'))$. Composing (CR.8) with the supplied bijection
from the level's double quotient to its finite class set $C_0$ gives
the literal map \texttt{clsR}, its equality criterion, and its
surjectivity. In particular,
\[
 \coprod_i\Gamma_i\backslash\mathrm{Lat}\ \simeq\ C_0.
                                                               \tag{CR.10}
\]

\paragraph{CR4. Why the actual stabilizers are trivial.}
First, every rational stabilizer of a level class is finite.
For an adelic representative $x$, intersect the transformed local
lattices $x_v\Lambda_v$ with $E^k$. Clearing the finitely many adelic
denominators sandwiches this intersection between
$a\mathcal O_E^k$ and $a^{-1}\mathcal O_E^k$ for a nonzero integer
$a$, so it is a full integer lattice. The stabilizer preserves it
and the Hermitian form, which is definite at every archimedean place.
In the Minkowski embedding, each of its basis vectors can therefore
have only finitely many possible images: a discrete lattice meets
a bounded set in finitely many points. Hence the stabilizer is finite.

At each of the source's two distinct auxiliary rational primes
$g_0,g_1$, a conjugate of a stabilizer element is integral and reduces
to an upper-unitriangular matrix. This is the literal
$\texttt{iwU}$ condition, weaker than being the identity modulo the
prime. We need the following elementary fact. If $A\in\mathrm{GL}_k(R)$
over a complete DVR of residue characteristic $g$ is upper
unitriangular modulo its maximal ideal $\mathfrak m$ and
$A^N=1$ with $(N,g)=1$, then $A=1$. Indeed its reduction has
$g$-power order, so it is already the identity. If
$A=1+X$ with $X\in M_k(\mathfrak m^r)$, then
\[
 A^N\equiv1+NX\pmod{\mathfrak m^{r+1}}
       \qquad(r\ge1).
\]
As $N$ is a unit, induction puts $X$ in every power of
$\mathfrak m$, whence $X=0$.

An element of finite order $N=g^a b$, $(b,g)=1$, in this local
subgroup consequently satisfies $A^{g^a}=1$. Thus its order is a
power of $g$. The maps from the rational field into either split
completion are injective; the global stabilizer element has both
$g_0$-power and $g_1$-power order, and is the identity.
Finally, if $\gamma\in\Gamma_i$ fixes $g\mathbf W$, then (CR.7)
puts $(y_i j(g))^{-1}\gamma(y_i j(g))$ in $U$.
The rational-stabilizer result therefore proves exactly the freeness
required in CR2, as well as triviality of the level-zero class
stabilizers.

\paragraph{CR5--CR6. Match the two dimensions.}
Combining (CR.5), (CR.6) and (CR.10) gives the general rank formula
\[
 \boxed{\ d=|C_0|\prod_{v,s}D_k(\lambda^0_{v,s}).\ }      \tag{CR.11}
\]
The supplied \texttt{Forms0} identifies the complex level-zero
automorphic forms with their coefficient values on $C_0$, subject
only to the class-stabilizer invariance conditions. CR4 removes
those conditions. Each coefficient fiber is the tensor of the
characteristic-zero Weyl representations of weights $\lambda_\tau$.
Consequently
\[
 \dim_{\mathbb C}\mathrm{Mu}_0
       =|C_0|\prod_{\tau\in\mathrm{Tau}}D_k(\lambda_\tau).
                                                               \tag{CR.12}
\]
The actual position map is
$\texttt{gatRec\_pos}=\texttt{eTau}^{-1}\circ\texttt{posG}$,
that is, $(v,s)\mapsto\texttt{eTau}^{-1}(\texttt{posG}(v,s))$;
it is a bijection. The definition \texttt{lamRec} identifies
$\lambda^0_{v,s}$ with precisely the corresponding $\lambda_\tau$.
The rank reindexing uses $m_{\mathrm{loc}}+1=k=N(0)$, while negative reversal at
odd positions was already accounted for in (CR.5). Thus
\[
 \boxed{\ (X.\texttt{gat}\ r\ e).d
          =\dim_{\mathbb C}(D.\texttt{mem}\ u).A.\mathrm{Mu}_0\ }
             \qquad(r,e\in\mathbb N).                    \tag{CR.13}
\]
This is the source's \texttt{locPack\_d\_eq}; substituting the
constructed \texttt{pack} gives \texttt{pack\_d\_eq}.

The exported wrapper takes \texttt{Joint0} as an argument but its
proof does not use it. This is mathematically exact: the two supplied
level descriptions already use the same finite set $C_0$ and the
same weights, which suffice for (CR.11)--(CR.12). Their actual
lattices, local subgroups and representatives need not be identified
to compare these dimensions. \texttt{Joint0} retains its separate
role in the character and comparison constructions. No integral
identification, Gram identity, metric transport, discarded-target
estimate, or global arithmetic window follows from (CR.13) alone.
